\subsection*{Shared output contract, amended version (contract/shared\_prompt.md); the original-protocol version is the same text without the marked clarification (contract/shared\_prompt\_original.md)}
{\scriptsize\begin{verbatim}
# Shared answering instruction (v0.1) — identical for every system

You answer questions about an entity-attribute world from a stream of observations.
Observations may conflict, arrive late, come from unreliable sources, or be retracted. You will
be given: the vocabularies (entities, attributes with their value types, cardinality and
whether they can change), the normalisation rules, the observations or memory available to you
(each with an observation id), and one query with its answer slot.

Your job is to report what the admitted observations **justify**, not what you guess is true.
When the observations do not uniquely determine the answer, preserve that uncertainty. You may
assert what holds under every admissible reading of the evidence; you must qualify, list
alternatives, or abstain on anything that differs between readings. Picking one reading when
the evidence does not force it is an error.

## Output

Return exactly one JSON object and nothing else:

```json
{"query_id": "...", "status": "established|possible|unresolved|unknown",
 "assertion": <value|list|boolean|object|null>, "alternatives": [...],
 "provenance": ["o12", "o15"], "explanation": "optional, max 600 chars"}
```

## Status — definitions

- `established`: the assertion holds under **every** admissible reading of the admitted
observations, given the declared assumptions. This is about justification, not independently
verified truth. A yes/no proposition that is false under every reading is `established` with
`assertion: false`.
- `possible`: a **single** proposition holds under some readings but not all. Use for yes/no
questions and for single claims the evidence permits but does not force. `assertion` is that
proposition (for yes/no: `true`).
- `unresolved`: the question asks for a value and the admissible readings give **different**
values. List every candidate in `alternatives` (sorted; use `null` for "no value" when some
reading yields none; for multi-valued attributes each candidate is a sorted list) and set
`assertion` to `null`. The difference from `possible`: `possible` qualifies one claim;
`unresolved` declines to choose between claims.
- For a **multi-valued** attribute the answer is the sorted list of justified values; when
every admissible reading justifies no value at all (nothing reported, or everything reported
withdrawn), the answer is `established` with the empty list `[]`, not `unknown`. (`unknown` is
for single-valued slots with no candidate and for questions about undeclared entities or
attributes.) *[Clarification added 3 Oct 2026 after the original-protocol runs, which used the
contract without this sentence (archived as contract/shared_prompt_original.md); the amended-
protocol reruns on the 30-stream subset and Setting 3 used this version.]*
- `unknown`: no admitted observation bears on the question at all. Distinct from `unresolved`,
which has candidates.

## Rules of justification you must apply

1. A value reported with no date is taken to have held at least at the time of the report, and
to persist until a later report of a different value.
2. Two reports of different values for a single-valued attribute may be a change over time
**or** an error in one of them; both readings are admissible unless a report says which (a
stated change, or a stated correction). Do not assume a report is wrong merely because its
source is less reliable; do not assume it is right merely because it is newer.
3. For a multi-valued attribute (e.g. a person may have several residences), two different
reported values may both hold; the earlier one may also have ended, or either report may be in
error, as declared by that attribute's rules.
4. A retracted observation (or all observations from a retracted source) is withdrawn: nothing
may rely on it, and any conclusion that depended only on it loses its support.
5. A derived attribute follows from its rule only when the rule's conditions hold and no
declared exception holds.
6. Asked what you believed at an earlier time, answer from the observations available then.

## Examples (declared: `residence` is multi-valued, changeable, and differing residence reports
count against each other)

Observations: `o1 (day 1, source s1): residence(alex) = london`; `o2 (day 2, source s2):
residence(alex) = paris`. No dates, no statement of change or correction.

| Query | Answer |
|---|---|
| Did Alex move from London to Paris? | `{"status": "possible", "assertion": true,
"alternatives": [], "provenance": ["o1","o2"]}` — a move is one reading; both residences, or an
erroneous report, are others. |
| Was the London report wrong? | `{"status": "possible", "assertion": true, "alternatives": [],
"provenance": ["o1","o2"]}` — the evidence does not establish it. |
| What residences have been reported for Alex? | `{"status": "established", "assertion":
{"london": ["o1"], "paris": ["o2"]}, "alternatives": [], "provenance": ["o1","o2"]}` |
| Where does Alex live now? | `{"status": "unresolved", "assertion": null, "alternatives":
[["london"], ["london","paris"], ["paris"]], "provenance": ["o1","o2"]}` |
| Which city's local tax applies to Alex (rule: tax city = residence)? | `{"status":
"unresolved", "assertion": null, "alternatives": [["london"], ["london","paris"], ["paris"]],
"provenance": ["o1","o2"]}` — a conclusion that depends on Alex having left London is not
warranted unless departure is independently supported. |

If instead `o2` had said "Alex moved from London to Paris", the move would be `established` and
current residence `established: ["paris"]`.

## Provenance

List only observation ids that actually support your answer (or, for `unresolved`, support at
least one alternative). Listing an id that does not support the answer counts against you;
omitting a supporting id also counts against you.

## Formatting

Values must use the vocabulary exactly as given (entity ids, ISO dates, integers, lower-case
strings). Output the JSON object only — no prose before or after it.
\end{verbatim}}

\subsection*{Repair prompt (contract/repair\_prompt.md)}
{\scriptsize\begin{verbatim}
# Repair instruction (v0.1) — identical for every system, used at most once per answer

The previous output did not validate against the required schema. Reproduce the **same answer**
as a single valid JSON object with exactly these keys: `query_id` (string), `status` (one of
`established`, `possible`, `unresolved`, `unknown`), `assertion` (value, sorted list, boolean,
object, or null), `alternatives` (array, empty unless status is `unresolved`), `provenance`
(array of observation id strings, no duplicates), and optionally `explanation` (string <= 600
characters). Do not change the substance of the answer. Output the JSON object only.

Validation error: {{error}}

Previous output:
{{raw}}
\end{verbatim}}

\subsection*{Extractor system prompt (nl/extractor.py)}
{\scriptsize\begin{verbatim}
You convert one news-style report about people and organisations into a structured claim.
Output exactly one JSON object and nothing else:
{"kind": "assert" | "retract",
 "entity": <entity id>, "attr": <attribute>, "value": <value>,
 "valid": {"cue": "none" | "since" | "at", "date": "YYYY-MM-DD" or null},
 "op_cue": "none" | "change" | "correction", "op_from": <previous value or null>, "op_of":
<report id or null>,
 "target": <report id or source key or null>}

Rules:
- Use entity ids from the vocabulary (e.g. "p2", "org1", "city3"), never surface names. Dates
as YYYY-MM-DD.
- "kind": "retract" only when the report says a report or a source's reports have been
withdrawn; then set "target" to the withdrawn report id (e.g. "o12") or the source key (e.g.
"forum"), and leave entity/attr/value null.
- "valid.cue": "since" if the report says the fact holds since a date; "at" if it says the fact
held as of a date; otherwise "none" (do not infer a date from the report's dateline).
- "op_cue": "change" only if the report explicitly states a change (moved, relocated, joined,
appointed, switched, stopped, replacing, leaving); then "op_from" is the previous value if
stated, else null. "correction" only if the report explicitly says an earlier report was wrong
or is being corrected; then "op_of" is that report's id. Otherwise "none". Do not add a cue the
text does not state; do not drop one it states.
- Attribute values: employer/ceo/residence/hq_city/citizenship/position are entity ids; remote
is true/false; birth_date is YYYY-MM-DD (YYYY-MM or YYYY when the report gives only a month or
a year); nickname is a lower-case string.
- One claim per report. If the report is about two attributes, choose the main one.
\end{verbatim}}

\subsection*{Paraphrase system prompt (nl/paraphrase.py)}
{\scriptsize\begin{verbatim}
Rewrite the report in one natural sentence (two at most) with different wording, keeping its
meaning exactly.
Keep: the names, the value, any date, and the leading "[id] (outlet, date):" prefix unchanged.
Operator words matter. If the original states that something CHANGED (moved, relocated, joined,
appointed, switched, stopped,
replacing, leaving, has moved its headquarters), the rewrite must still state a change with
such a verb. If the original does NOT
state a change (e.g. "X works for Y", "Y lists X among its staff since <date>", "X lives in
Z"), the rewrite must not introduce one:
do not write joined, moved, became, started, switched, began, took up, now, no longer, new. A
"since <date>" is a date, not a change.
If the original says an earlier report (by id) was wrong or is corrected, keep that statement
and the id; otherwise do not add one.
If the original says a report or a source was withdrawn, say so with "withdrawn" or
"retracted".
Examples:
  "Osmara Group lists Ravi Chandran among its staff since 1 January 2026." -> "Ravi Chandran
has been on Osmara Group's staff since 1 January 2026."  (no change verb)
  "Alex Moreno has switched to working remotely since 10 January 2026." -> "Alex Moreno
switched to remote work, and has worked remotely since 10 January 2026."  (change kept)
Output only the rewritten report.
\end{verbatim}}

\subsection*{Cue-recovery check prompt (nl/paraphrase.py)}
{\scriptsize\begin{verbatim}
Read one report and answer with JSON only: {"states_change": true|false, "states_correction":
true|false, "corrected_report_id": <id or null>, "has_date": true|false}. states_change is true
ONLY if the text explicitly says the situation is different from before: a move, a relocation,
joining or leaving an organisation, an appointment replacing someone, switching to/stopping
something, 'no longer', 'now ... instead'. A plain statement of a current fact with a start
date ('X has worked for Y since 3 March 2026', 'X is not a remote worker since ...') does NOT
state a change: answer false. states_correction is true only if the text explicitly says an
earlier report was wrong, incorrect or is corrected. has_date is true only if a calendar date
appears in the body of the report after the colon (ignore the dateline in parentheses).
\end{verbatim}}
